\section{Uniqueness of the limiting background paths}
\label{sec:background-path-convergence}
The path coefficients require their own convergence argument. Uniform
bounds give convergent subsequences, but do not identify their limits.
The identification uses the order of the Taylor coefficients in the
stationary action: first the cubic term, then the quartic term, and
finally the quadratic loop correction. We explain both the finite-step
comparison and the passage from subsequences to the whole sequence.

At each scale $m$ write a path as its prescribed one-form row and three
coefficient arrays
\[
                    (a_m,b_m,c_m).
\]
Here $a_m$ represents the quadratic loop correction, $b_m$ the cubic
term, and $c_m$ the quartic term. Each array belongs to the selected
closed subspace of its real coefficient Hilbert space. The selected
subspace is the orthogonal complement of the kernel of all finite-period
Taylor readouts. Consequently two selected arrays with the same
readouts on every prescribed period are equal. This statement concerns
the arrays themselves, not just their evaluation on one torus.

For $q=2,3,4$ and the fixed encoding width $w_{m,q}$, use the norm
\[
 \mathfrak n_{m,q}(C)
   =\sum_{a,\nu}e^{-(w_{m,q}-v_m)\ell(a)}|C(a,\nu)|,
 \qquad w_{m,q}>v_m.
\]
Here $v_m$ is the paid path width, and the canonical lengths and finite
color index $\nu$ are as in Section~\ref{sec:fixed-rung-convergence}.
All weights are strictly positive. Thus this expression is a norm on
the coefficient space on which it is finite; homogeneity and the
triangle inequality follow term by term. Let $s_{m,q}>0$ be its paid
allowance. The ordinary outputs satisfy
$\mathfrak n_{m,q}\le s_{m,q}$. The extended input domain permits
$64s_{m+1,q}$, and keeps the other two arrays at their ordinary
allowances. The same finite-step construction is available for every
input in this domain. We use this actual extended domain below;
existence for one selected input alone would not suffice.

\subsection{The triangular affine identities}
Fix the fine and coarse one-form rows. Consider three admissible
finite-step inputs and outputs, indexed by $0,1,t$, with the same
rows and the same forest, coordinates and normalization conventions.
The following identities hold whenever the indicated interpolation
is admissible:
\begin{align}
 b^{\rm in}_t=b^{\rm in}_0+t(b^{\rm in}_1-b^{\rm in}_0)
 &\ \Longrightarrow\
 b^{\rm out}_t=b^{\rm out}_0+t(b^{\rm out}_1-b^{\rm out}_0),
 \label{eq:background-cubic-affine}\\
 c^{\rm in}_t=c^{\rm in}_0+t(c^{\rm in}_1-c^{\rm in}_0)
 &\ \Longrightarrow\
 c^{\rm out}_t=c^{\rm out}_0+t(c^{\rm out}_1-c^{\rm out}_0),
 \label{eq:background-quartic-affine}\\
 a^{\rm in}_t=a^{\rm in}_0+t(a^{\rm in}_1-a^{\rm in}_0)
 &\ \Longrightarrow\
 \begin{cases}
 a^{\rm out}_t=a^{\rm out}_0+t(a^{\rm out}_1-a^{\rm out}_0),\\
 \alpha_t=\alpha_0+t(\alpha_1-\alpha_0).
 \end{cases}
 \label{eq:background-loop-affine}
\end{align}
In \eqref{eq:background-quartic-affine} the three input cubic arrays
must agree. In \eqref{eq:background-loop-affine} both their cubic
and quartic arrays must agree. No equality of those arrays is needed
in \eqref{eq:background-cubic-affine}.

Here is the stationary-action argument for the first two identities.
In finite-dimensional chart coordinates let $S(x,y)$ be the analytic
action, $y=\psi(x)$ its stationary graph, and
$S_{yy}(0,0)=H$ its invertible Hessian in the eliminated coordinates.
All these stationary graphs have $\psi(0)=0$. Write $z(x)=(x,\psi(x))$.
Differentiating $S_y(x,\psi(x))=0$ once determines $D\psi(0)$
from the common quadratic jet of $S$. Repeated differentiation
determines each next derivative by an equation
\[
 H D^r\psi(0)=-R_r,
\]
where $R_r$ is a finite sum of contractions of derivatives of $S$
of order at most $r+1$ and derivatives of $\psi$ of order less
than $r$. This follows directly by the multivariable chain rule:
the only term containing $D^r\psi$ is obtained by differentiating
the $y$ argument of $S_y$ once and the graph $r$ times.

Suppose three actions have the same jet through order $r$, with
$r=2$ or $3$, and their order-$(r+1)$ jets have the displayed
affine relation. Their graph jets through order $r-1$ agree.
In $D^{r+1}(S\circ z)(0)$, terms containing a single highest
graph derivative cancel by stationarity. Terms containing
$D^r z$ pair with $D^2S$ and $Dz$ and vanish by the derivative
of stationarity, since $D^r z$ is vertical. All remaining terms
involve the common lower jets, except the single term
\[
                       D^{r+1}S(0)(Dz(0),\ldots,Dz(0)).
\]
It follows that the reduced order-$(r+1)$ jet has the same affine
relation. Equivalently, one may subtract the three chain-rule
expansions: every common lower-jet term cancels, and the displayed
term is the only possible difference.

Cubic path terms vanish through order two, and quartic terms
through order three. The quadratic action depends only on the
common row. Thus the preceding calculation with $r=2$ proves
the affine cubic identity for the reduced actions. The finite-step
classical matching equates these jets with the coarse path jets.
Subtracting their common coarse-row contribution and using the
injectivity of the selected cubic readout gives
\eqref{eq:background-cubic-affine}. If the input cubic arrays
agree, this identity first makes the output cubic arrays agree.
The calculation with $r=3$ then applies; subtracting the common
coarse-row and cubic contributions gives
\eqref{eq:background-quartic-affine} by the quartic readout.

For the loop identity the classical action, including its cubic
and quartic terms, is the same for all three inputs. The implicit
function theorem and the invertibility of $H$ give the same
stationary germ. Its Hessian determinant and the coordinate
Jacobian contribution are therefore the same. The remaining
one-loop expression is the loop correction evaluated on that
germ; it is linear in $a^{\rm in}$. Its quadratic jet is affine.
The coarse one-loop decomposition has the form
\[
                    \alpha Q_0+Q_a+\mathcal E,
\]
where $Q_0$ is the common normalized quadratic row form, $Q_a$
is the loop readout, and $\mathcal E$ is the finite-period error.
The latter has the proved exponential bound in the period.
Subtract the three decompositions. The all-period detection
statement for the quadratic forms makes the remaining cover
form zero. Its normalization functional takes the value one on
$Q_0$ and zero on the selected loop subspace. Applying it gives
the identity for $\alpha$. The residual loop readouts then vanish
on every period; injectivity on the selected loop subspace gives
the identity for $a^{\rm out}$. This identifies exactly where
quadratic normalization and all-period detection enter the proof.

\subsection{An enlarged input ball gives contraction}
Choose one of the three arrays, subject to equality of its preceding
arrays in the order $b,c,a$. Write $s'=s_{m+1,q}$ and $s=s_{m,q}$.
Let $x,y$ be two ordinary inputs and put
$z=\mathfrak n_{m+1,q}(y-x)$. If $z=0$, then $x=y$, and the
affine identity at parameter zero gives equal selected outputs.
If $z>0$, use the real parameter
\[
                        t=32s'/z.
\]
The replacement $x+t(y-x)$ has norm at most
\[
             \mathfrak n_{m+1,q}(x)+tz\le33s'\le64s'.
\]
Keeping the other two arrays fixed preserves their invariant
subspaces and bounds. In the loop case it also preserves the
linear loop normalization constraints. Thus this is an allowed
extended input. Its output and that of $x$ both have norm at
most $s$. By the affine identity, their difference is $t$ times
the difference of the original two outputs. Hence
\[
 t\mathfrak n_{m,q}(y^{\rm out}-x^{\rm out})\le2s,
 \qquad
 \frac{\mathfrak n_{m,q}(y^{\rm out}-x^{\rm out})}{s}
       \le\frac1{16}\frac{z}{s'}
       \le\frac14\frac{z}{s'}.
\tag{$\ast$}\label{eq:background-contraction}
\]
We retain the weaker factor $1/4$ used in the history argument.
The extrapolated parameter need not lie in $[0,1]$; the extended
input domain is essential for this proof.

\subsection{Passing finite histories to cluster histories}
Fix any sequence of terminal cutoffs tending to infinity. At each
fixed scale, the three coefficient arrays have uniformly bounded
Hilbert norms and the scalar corrections $\alpha$ are bounded.
Enumerate the countably many scales and coefficient coordinates.
Successive extraction in bounded real intervals, followed by the
diagonal subsequence, gives convergence of all those coordinates
and all scalars at once. Finite sums of coordinate squares show
that each limiting array has the same Hilbert bound. A closed
linear constraint is retained: for a test vector in its orthogonal
complement, approximate by a finite-coordinate vector; the uniform
Hilbert bound controls the omitted inner product. Thus the limits
remain in the selected subspaces.

The one-form rows have their common limit from the retained kernel
horizon. The width argument in
Section~\ref{sec:fixed-rung-convergence} gives convergence of
the path sums at every smaller width. It also gives convergence
of every finite chart derivative. Indeed, on a sufficiently
small complex chart ball each length-$\ell$ path and each fixed
derivative costs a polynomial in $\ell$ times $e^{\epsilon\ell}$.
Choose $\epsilon$ smaller than the available width difference;
the remaining exponential factor sums the polynomial. Termwise
differentiation is then justified by uniform absolute convergence.

The limiting eliminated Hessian obeys the same positive lower
bound, by passage to the limit in its quadratic form. The analytic
implicit function theorem supplies a stationary germ for the
limiting action. The derivative recursion above proves convergence
of every derivative of the finite stationary germs: inversion
of the finite Hessian is continuous on the matrices with this
common lower bound, and each $R_r$ is a finite expression in
previously convergent derivatives. The limits also have the same
classical matching through degree four. For the one-loop expression,
the derivatives of $\log\det H$ are finite sums of products of
$H^{-1}$ and derivatives of $H$. The positive lower bound and
the same recursion therefore give convergence of its quadratic
jet, including the fixed coordinate Jacobian. All finite-period
error inequalities pass to the limit. The row normalization does
too. Thus the resulting cluster history satisfies the same local
relations and paid output bounds as the finite histories.

The input at scale $m$ equals the output at scale $m+1$ in these
cluster histories: it is an exact finite-history identity on
every sufficiently late cutoff, so coordinatewise passage to
the limit preserves it. All periods and forests needed here
have already been fixed by the schedule. The derivative recursion
defines their limiting stationary germs from the same limiting
path; no independent subsequence is chosen for each germ.

\subsection{Comparison of two cluster histories}
Let $F$ and $H$ be two cluster histories with the same retained
row horizon. A finite-step comparison cannot simply be applied
to them as if they came from the same finite cutoff. Instead fix
a scale and an array type, with its preceding input arrays equal.
Take finite steps tending to $F$, and replace just this array in
each fine input by the corresponding array of $H$. The replacement
is an ordinary admissible input: the bounds and linear invariance
constraints hold separately for its three arrays. Construct its
output at that very same cutoff. Extract a further subsequence
for these modified outputs and their scalar corrections.

The preceding limit argument gives a valid cluster step. Its
chosen input array is that of $H$, and its preceding arrays are
also those of $H$. The affine identity at parameter zero shows
that its chosen output array equals that of $H$. Apply
\eqref{eq:background-contraction} before taking the subsequence
limit. The weighted norms of the differences converge by the
strict-width dominated convergence argument. We obtain
\[
 g_m\le\tfrac14 g_{m+1},\qquad
 g_m:=\frac{\mathfrak n_{m,q}(H_m-F_m)}{s_{m,q}},
                  \qquad 0\le g_m\le2.
\]
Iteration gives $g_m\le2\,4^{-h}$ for every integer $h\ge0$;
hence $g_m=0$. Apply this first to the cubic arrays at every
scale, which require no preceding equality. Then apply it to
the quartic arrays, using the cubic equality just obtained.
Finally apply it to the loop arrays, using both equalities.
All three arrays in the two cluster histories agree.

Their rows already agree, so their full input and output paths
agree. The loop identity at parameter zero identifies their
scalar corrections. Uniqueness of the stationary germ identifies
all its derivatives for every fixed forest and period. We have
therefore proved uniqueness of the entire cluster history,
including the scalar and stationary data needed in the next
scale comparison.

If the original sequence failed to converge to this history in
any one coordinate, scalar, or fixed chart derivative, a
subsequence would stay outside some neighborhood of its asserted
limit. The compactness construction applies to that subsequence
and gives a further subsequence with the unique cluster limit,
a contradiction. Thus the entire cutoff sequence converges.
The strict-width argument upgrades its coefficientwise convergence
to the weighted path convergence used for fixed-scale normalized
moments. This closes that part of the fixed-scale input once the
finite-step extended-domain construction, retained row horizon,
and canonical all-period detection used above are supplied on
the selected history.
